\chapter{Discusión}
\label{cap:discusion}

Este capítulo interpreta las medidas del capítulo «\nameref{cap:evaluacion}», las sitúa respecto a otros enfoques de la
literatura, recoge las limitaciones del trabajo y discute qué hace falta para aplicar el método a otras redes.

\section{Interpretación de los Resultados}
\label{sec:interpretacion}

Los resultados permiten distinguir de dónde sale la precisión del método.

\subsection{El Tramo en que Está el Autobús} Saber en qué tramo está el autobús aporta la mayor parte de la información. La capa 2, que solo conoce la
próxima parada y la hora del dato, queda a pocos puntos de la capa 1 con GPS (Tabla~\ref{tab:capas}). La posición
exacta dentro del tramo mejora la predicción, pero lo decisivo es partir de dónde está el autobús y no de la hora
publicada, como anticipa la proposición~\ref{prop:diferencias}.

\subsection{La Corrección del Horario} Corregir el error de forma del horario ya lo mejora mucho. La capa 3 no sabe
nada del autobús y, aun así, duplica o triplica los aciertos del horario publicado. Esa mejora solo puede venir de
sustituir las duraciones interpoladas por los tiempos aprendidos de cada tramo. Le quedaba, sin embargo, un sesgo:
llegaba de media 131~s antes que el autobús. La corrección por viaje, que usa lo que cada viaje suele hacer y nada del
autobús ese día, lo deja en 17~s (sección «\nameref{sec:eval-correccion}»). Quitar el sesgo de la media no es lo mismo
que acertar lo típico: con el percentil 40 de los desvíos en lugar de la media, la hora sin autobús acierta más y el
autobús pasa antes de lo mostrado un tercio menos de veces (sección «\nameref{sec:correccion-percentil}»). Lo que
queda es el desvío de cada viaje ese día, que se explica por el autobús que lo hace. Si ese desvío se conociera, el
acierto a $\pm2$~min de las filas «Estimado» sin autobús del 2 y el 3 de noviembre pasaría del 52--58\,\% (con la
media) al 81--85\,\%; pero ni los demás viajes de la línea ni la hora que ya ha pasado sin autobús lo predicen.

\subsection{La Cobertura de Tramos Aprendidos} La cobertura de tramos aprendidos es lo que más mueve la precisión. En
las llegadas a menos de diez minutos, con todo el camino aprendido se acierta el 93\,\% dentro de $\pm2$~minutos,
frente al 76\,\% con un tramo sin aprender y el 59\,\% con dos o más (Tabla~\ref{tab:cobertura} y el texto que la sigue; con 17 días, el 93, el 76 y el 63\,\%,
Tabla~\ref{tab:varios-dias-origen}). La curva de aprendizaje (Tabla~\ref{tab:curva}) indica que esa
cobertura sigue creciendo con los datos.

\subsection{El Error y la Distancia} El error crece con la distancia, y el acierto tiene un techo. Cada tramo suma su
propia variación, de modo que el
acierto baja de forma regular con el horizonte. En las llegadas a menos de diez minutos con todo el camino aprendido, el techo medido
está en torno al 93\,\% dentro de $\pm2$~minutos, limitado por la resolución del dato (una posición por minuto y $\pm30$~s en el paso real) y por lo que
ningún modelo puede prever con estos datos: semáforos, tráfico y esperas en parada.

\section{Relación con Otros Enfoques}
\label{sec:otros-enfoques}

El método pertenece a la familia de las medias históricas (capítulo «\nameref{cap:marco}»), con dos
diferencias: los tiempos se aprenden por tramo y franja en lugar de por línea, y la posición en tiempo real reinicia la
predicción en cada actualización. Otros trabajos han usado filtros de Kalman con datos de localización y de conteo de
pasajeros \citep{shalaby2004} o redes neuronales \citep{chien2002,jeong2005}; en el caso de \citet{jeong2005}, las
redes neuronales superaron a un modelo de datos históricos y a modelos de regresión en una línea de Houston.

Las cifras de este trabajo no se pueden comparar directamente con las de esos estudios: las redes, los datos, los
horizontes y las medidas de error son distintos. Por eso se implementaron esos métodos sobre las mismas muestras y con
el mismo protocolo por días (sección «\nameref{sec:comparacion}»). El resultado tiene dos lecturas. La primera es que
el modelo sencillo se queda con casi toda la mejora sobre el horario: de 154 a 42~s de error medio a 4--8~minutos,
mientras que la mejor alternativa, una red neuronal que además usa lo que pasa ese día, baja a 39,5~s. La segunda es
que esa diferencia, aunque pequeña, es real: la red reduce el error medio entre un 6 y un 7\,\% y mejora al método
todos los días de prueba en todos los horizontes, y un filtro de Kalman sobre el ritmo del propio autobús consigue algo más de la mitad de esa mejora con dos
parámetros.

Esto no contradice que corregir con la velocidad reciente del autobús empeorara la medida
(sección «\nameref{sec:descartadas}»). Esa corrección usaba la velocidad del último minuto, muy ruidosa, sobre la llegada
completa; el filtro usa la desviación de los tramos ya terminados, la suaviza y la olvida con una constante de 30~minutos.
Que los autobuses previos en el mismo tramo aporten tan poco (un 1 o 2\,\%) indica que, en esta red, la variación de
un día a otro en un tramo pesa menos que la de cada conductor y cada viaje, que es la que capta el ritmo del propio
autobús. La comparación se hizo sobre sumas de tramos, no sobre la llegada completa que ve el usuario; antes de
adoptar el filtro habría que medirlo así, con la regla de la sección «\nameref{sec:eval-contraccion}».

\section{Limitaciones}
\label{sec:limitaciones}

\begin{itemize}
  \item \textbf{Una muestra pequeña con las tres capas.} Las Tablas~\ref{tab:resultados} a~\ref{tab:cobertura}, las
        únicas que calculan cada llegada con las tres capas a la vez, salen de 14~minutos de feed de un laborable; sus cifras a horizonte
        largo tienen intervalos de confianza anchos. La medida principal se repitió con 17 días
        (Tabla~\ref{tab:varios-dias}), y con ellos el acierto a 10--20~minutos quedó en el 67\,\% y no en el 73\,\%, ya
        con llegadas de cualquier capa.
  \item \textbf{Un método que cambia.} Los 17 días miden la aplicación tal como estaba cada día, con calibraciones y
        reglas que cambiaron en ese periodo; no son la precisión de una versión fija del método.
  \item \textbf{Resolución de la verdad de referencia.} El paso real se conoce con $\pm30$~s. Las medidas dentro de
        $\pm1$~minuto están especialmente afectadas.
  \item \textbf{Definición de paso.} La verdad de referencia por cambio de parada mide cuándo el feed cambia de
        próxima parada, lo que puede ir algo por detrás del paso físico. Afecta por igual al método y al horario publicado.
  \item \textbf{Censura.} Solo se evalúan los autobuses cuyo paso se observa dentro de la ventana de captura; los muy
        retrasados pueden quedar fuera.
  \item \textbf{Dependencia temporal.} Las predicciones de un mismo autobús en instantáneas sucesivas están
        correlacionadas; el número efectivo de observaciones independientes es menor que el número de predicciones.
  \item \textbf{Temporada y festivos.} Las muestras van del 6 de octubre a los primeros días de noviembre de 2026 y
        solo tienen un festivo entre semana. La regla de
        los festivos se apoya en ese día, y el efecto de la temporada (verano, vacaciones) está sin medir.
  \item \textbf{Lo que el feed no dice.} Un viaje sin autobús en el feed sale con la capa 3 aunque no se haga, y un
        refuerzo que no está en el horario no sale en las paradas. Según las capturas del 3 de noviembre de 2026, en
        torno al 4\,\% de los viajes que deberían estar en marcha no se pueden interpretar: los que salen tarde y los
        que no aparecen en todo el día se ven igual en el momento (sección «\nameref{sec:feed}»).
  \item \textbf{De madrugada.} De 23:58 a 05:00 el feed de Bizkaibus repite la lectura de las 23:58 y los autobuses
        nocturnos salen con la capa~3 o con el horario (sección «\nameref{sec:mas-de-24h}»). Esas llegadas no están
        medidas, y la parte de un viaje posterior a la medianoche no aporta muestras al aprendizaje.
  \item \textbf{Calidad del horario.} El método depende de que trazado y paradas casen. En el horario de Bizkaibus del
        3 de noviembre de 2026, 1\,417 de los 6\,130 viajes (el 23\,\%) tenían alguna parada a más de 200~m del
        punto del trazado que marcan sus metros, 670 de ellos (el 11\,\% del total) con el trazado al revés. La
        corrección de la sección «\nameref{sec:trazados-gtfs}» lo arregla en la construcción, pero otra red puede
        traer errores distintos, que habría que buscar antes de medir el método.
  \item \textbf{Una sola red.} Todos los resultados son de Bizkaibus. La generalización a otras redes es una hipótesis
        razonable (sección «\nameref{sec:aplicabilidad}»), no un resultado.
  \item \textbf{Evaluación hecha por el autor.} El método lo ha evaluado quien lo ha diseñado, con sus propias
        herramientas. Los vectores de prueba y la implementación de referencia permiten comprobar que el método descrito
        es el que se ha medido, pero no sustituyen a una replicación independiente con otros datos.
  \item \textbf{Sin revisión por pares.} Este documento no ha pasado por una revisión externa.
\end{itemize}

\section{Límites de la Infraestructura}
\label{sec:infraestructura}

Todo el sistema funciona sin un servidor propio y casi sin coste, y eso condiciona el método tanto como los datos.

\begin{itemize}
  \item \textbf{La aplicación, sin estado.} La aplicación y la API corren en una plataforma sin servidor (Vercel): cada
        petición es una función independiente que no recuerda nada de la anterior, salvo una caché corta en disco y un
        almacén clave-valor externo (Upstash Redis) que se paga por comando. Cada función admite 250~MB, de los que la
        aplicación usa 208, y cada red nueva suma su base de datos del horario.
  \item \textbf{El aprendizaje, en una plataforma de pruebas.} La captura y la recalibración corren en una plataforma
        de integración continua (GitHub Actions) pensada para compilar y probar, no para trabajar sin parar: cada trabajo
        dura como mucho seis horas y los disparos programados se retrasan o se saltan (sección «\nameref{sec:ciclo}»).
        Hasta el 31 de octubre de 2026 solo se ejecutaron 163 de unas 600 capturas previstas.
  \item \textbf{Los datos, en ficheros.} Las muestras acumuladas no están en una base de datos, sino en ficheros JSON
        Lines guardados como adjuntos del repositorio. Cada recalibración los descarga enteros, los recalcula en memoria
        y los vuelve a subir. Por eso existen los topes de la sección «\nameref{sec:retencion}»: 120~días, 300 muestras
        por tramo y tipo de día, y 500\,000 muestras de ritmo, con ficheros de unos 100--120~MB y menos de 2~GB de memoria.
  \item \textbf{Una vez al día.} Los parámetros se recalculan una vez al día y producción usa parámetros de entre 23 y
        47~horas de antigüedad: un desvío o unas obras de hoy no entran en lo aprendido hasta uno o dos días después.
\end{itemize}

\noindent Estos límites explican varias limitaciones del método:

\begin{itemize}
  \item \textbf{Memoria corta.} Con 120~días y un tope por tramo no hay un año de historia. No se puede aprender una
        fiesta con la misma fiesta del año anterior (Navidad, Semana Santa, Aste Nagusia), ni el verano o las
        vacaciones escolares, y los días especiales se mezclan con los normales al aprender.
  \item \textbf{Grupos gruesos.} Con pocas muestras por tramo, lo aprendido se agrupa en seis franjas y dos tipos de
        día. Con más historia se podría separar por día de la semana, por hora o por fecha especial sin quedarse sin
        datos en cada grupo.
  \item \textbf{Sin otras fuentes.} No se cruzan la meteorología, el calendario escolar, los eventos (un partido o un
        concierto) ni las incidencias publicadas, porque no hay dónde guardar su historia junto a la de los pasos.
  \item \textbf{Sin modelos más capaces.} La red neuronal de la sección «\nameref{sec:comparacion}», que usa además
        lo que está pasando ese día, reduce el error medio de la suma de tramos entre un 6 y un 7\,\% frente al método. Se descartó para producción porque lo
        sencillo se puede auditar, pero también porque entrenarla y servirla en directo exige un servidor que no existe.
\end{itemize}

\noindent Con inversión en infraestructura se podría quitar cada uno de estos límites:

\begin{itemize}
  \item un servidor que capture el feed sin pausa, sin huecos ni capturas saltadas;
  \item una base de datos de series temporales que guarde años de pasos sin tope y permita consultar solo lo necesario,
        en lugar de descargar y recalcular ficheros enteros;
  \item aprender por fecha especial (festivos, Navidad, fiestas de cada pueblo, vacaciones escolares) comparando con
        el mismo día del año anterior, y cruzar los pasos con la meteorología, los eventos y las incidencias;
  \item recalibrar cada hora o de forma continua, de modo que un desvío entre en lo aprendido en horas y no en días;
  \item servir en directo modelos más capaces, como la red neuronal ya medida, e intervalos de predicción;
  \item añadir redes sin el límite de tamaño de cada función.
\end{itemize}

\noindent Ninguna de estas mejoras está medida. Son la dirección que señalan las medidas de este trabajo: la cobertura
de tramos aprendidos es lo que más mueve la precisión (sección «\nameref{sec:cobertura-tramos}»), y más historia y más fuentes
dan más cobertura y grupos más finos. Cuánto mejoraría la precisión en festivos o en verano solo se sabrá con esos
datos y con el mismo protocolo de medida.

\section{Aplicabilidad a Otras Redes}
\label{sec:aplicabilidad}

El método se diseñó para Bizkaibus, pero casi nada de él es específico de esa red. Esta sección separa lo que es
general de lo que hay que adaptar y da una lista de pasos para poner el método en marcha en otra ciudad. Describe
además las otras dos situaciones que ya resuelve la aplicación: redes que publican su propio retraso, que se tratan en
el motor, y redes que publican su propia predicción, que se tratan fuera de él.

\subsection{Requisitos}
\label{sec:requisitos}

El método necesita, por orden de importancia:

\begin{enumerate}
  \item \textbf{Un horario GTFS} con viajes, paradas y horas de paso. Sin él no hay recorrido sobre el que situar al
        vehículo.
  \item \textbf{Un feed en tiempo real que identifique el viaje} de cada vehículo y la \textbf{próxima parada}, con la
        hora del dato. Con eso funcionan las capas 2 y 3 y el aprendizaje queda limitado a lo que se pueda reconstruir.
  \item \textbf{La posición GPS} de cada vehículo y la hora en que se tomó. Con ella funcionan la capa 1 y el
        aprendizaje.
  \item \textbf{Trazados} (\texttt{shapes.txt}). Si no traen los metros de cada parada
        (\texttt{shape\_dist\_traveled}), la construcción de la base los mide sobre el trazado; sin trazados, la
        localización usa rectas entre paradas, con un pasillo más ancho, y no hay respaldo por longitud.
  \item \textbf{Una forma de capturar} el feed periódicamente y de guardar las muestras: aquí, un trabajo
        programado de una plataforma de integración continua y un fichero adjunto a una publicación del repositorio.
\end{enumerate}

\noindent Un operador que no publique retrasos ni predicciones, como Bizkaibus, es precisamente el caso para el que el
método aporta más. Uno que sí los publique puede usarlos directamente (sección «\nameref{sec:otras-redes}»).

\subsection{Qué Hay que Adaptar}
\label{sec:adaptar}

\subsubsection{La Lectura del Feed y el Emparejamiento} Es la única parte que depende del formato concreto del operador:
cómo se extraen la línea, el número de viaje y la salida de cabecera de la referencia del viaje
(expresión~\eqref{eq:regex}) y cómo se encuentra el viaje del GTFS (sección «\nameref{sec:emparejamiento}»). Si la referencia
del feed coincide con el \texttt{trip\_id} del GTFS, basta la tabla de alias. El motor deduce además de qué día es cada
viaje por la salida de cabecera del feed~\eqref{eq:dia-viaje}, que en Bizkaibus va en horas del horario. Los lectores
que la toman de \texttt{OriginAimedDepartureTime}, en hora del reloj (Bilbobus, Lurraldebus), considerarían de hoy un
viaje que empezara pasadas las 24:00; en esas dos redes no hay ninguno.

\subsubsection{Las Franjas} Las seis franjas por tipo de día responden a la demanda de Bizkaia. Una ciudad con otra
distribución de la hora punta debería revisarlas, con el mismo protocolo de evaluación por días.

\subsubsection{Las Constantes Medidas} Las constantes del Apéndice~\ref{ap:constantes} son de dos clases. Unas son
geométricas o de diseño y no deberían cambiar (la ventana de tramos y la regla de desempate). Otras se midieron en Bizkaibus y conviene volver a medirlas: el pasillo de 50~m depende de la calidad de los
receptores GPS, los 20~minutos de la capa 2 dependen de la longitud típica de los tramos, el tope de 45~minutos de los
retrasos reales de la red y los 60~s de salida de cabecera de la práctica de los conductores. La mitad del tramo de la capa 2 también se midió:
fue el mejor reparto entre 0,35 y 0,65.

\subsubsection{Sin Muestras} En una red nueva, sin muestras propias (la semilla del repositorio solo sirve para las
paradas de Bizkaibus), la tabla de tiempos por tramo está vacía y el modelo se reduce al factor de
ritmo, que sin muestras vale 1: las diferencias de horario de la proposición~\ref{prop:diferencias}, sin corrección de
forma. Solo con eso ya mejora mucho al horario publicado, porque elimina el retraso acumulado. Con unos días de capturas aparecen los
primeros tramos aprendidos y el modelo mejora sin intervención.

\subsection{El Mismo Motor en Otras Redes de Autobús}
\label{sec:lurraldebus}

El motor se reutiliza por herencia en dos redes más. Lurraldebus (Gipuzkoa) recibe su información en tiempo real en
varios feeds SIRI-VM, uno por operador; Bilbobus (los autobuses urbanos de Bilbao) la recibe en un solo feed del
Consorcio de Transportes de Bizkaia. En las dos, una clase derivada del emparejador solo cambia la red a la que se
aplica y los ficheros con los datos aprendidos (factor de ritmo, tiempos por tramo y correcciones por viaje); el
cálculo es el mismo código. La lectura del feed sí es distinta: un lector genérico toma la referencia del viaje tal
cual la publica el operador y la empareja con los alias del horario, con el nombre del operador delante cuando hay
varios. Las dos redes están activadas, pero sus feeds llegan vacíos durante largos periodos: entonces la aplicación
muestra el horario oficial sin afirmar que sigue a ningún vehículo, y el aprendizaje de tramos no tiene con qué
empezar hasta que los feeds traigan vehículos de forma regular. Sus GTFS no traen los metros de cada parada; la
construcción los mide sobre el trazado (Tabla~\ref{tab:sin-metros}), así que, cuando lleguen vehículos, se situarán
con el mismo pasillo de 50~m que Bizkaibus y no con la recta entre paradas.

\subsection{Redes que Publican Retraso o Predicción}
\label{sec:otras-redes}

\subsubsection*{Retraso Declarado}

En el código, «red de autobús» es la red para la que se configuró el emparejador; las demás, sean o no de autobuses,
usan este camino. En ellas las capas 2 y 3 y el contraste no se aplican, y la entrada se considera en directo si no tiene
posición o si la tiene y es válida. La llegada se calcula a partir del retraso declarado por el operador, $\delta_{\mathrm{decl}}$:
\begin{equation}
\ETA_{\mathrm{plano}}=a_\dest+\delta_{\mathrm{decl}}.
\end{equation}
Si la entrada no tiene posición pero sí parada actual y orden $n$, y el horario da la llegada $a_n$ a esa parada, se
calcula también una llegada por diferencias desde ahora,
\[
\ETA_{\mathrm{pos}}=\tnow+(a_\dest-a_n),
\]
y se usa en lugar de la plana si $a_\dest-a_n\ge0$ y $|\ETA_{\mathrm{pos}}-\ETA_{\mathrm{plano}}|\le900$~s
(\texttt{POSITION\_SANITY\_SECONDS}). Con posición válida se usa la capa 1 tal cual. En estas redes no se aplica el
tope de 45~minutos ni el estado \texttt{estimated}.

\subsubsection*{Predicciones de Llegada (GTFS-Realtime Trip Updates)}

Metro Bilbao, Euskotren, los tranvías de Bilbao y Vitoria-Gasteiz y Álavabus publican, para cada vehículo, la hora
prevista en cada parada que le queda; Renfe, Tuvisa y Dbus, un retraso por viaje. Ahí no hace falta estimar: basta con
emparejar cada vehículo con su viaje del horario y comprobar que lo que dice el operador es un dato en directo. Ese
trabajo no aprende nada, así que está en la aplicación y no en el motor.

\paragraph{Emparejar por el Viaje} Todos esos feeds traen el identificador del viaje del horario (con el horario del
4 de noviembre de 2026: los 71 viajes de Metro en servicio, 188 de Euskotren, 46 y 73 de los tranvías), y se empareja
por él. La hora en una parada es la del feed; si el feed se salta esa parada, la de la siguiente que trae, menos lo que
se tarda hasta ella según el horario; con un retraso por viaje, la del horario más ese retraso. El retraso que se
muestra es la diferencia con la hora publicada, redondeado a cero por debajo de 60~s (\texttt{ON\_TIME\_SECONDS}).
Las horas del feed son del reloj y las del horario pueden pasar de las 24:00: cada viaje guarda los días enteros que
separan unas de otras (24~h para un tren anotado a las 26:00 que pasa a las 02:00; 0 el resto del día), y todas las
comparaciones los descuentan. Sin eso, ese tren salía con $\pm1\,440$~min de retraso.

La versión anterior emparejaba por paradas y hora: cada tren del feed con la fila del horario más cercana dentro de
10~minutos, asignando de menor a mayor diferencia. Se suponía que con trenes separados varios minutos la ambigüedad era
rara. El 4 de noviembre, con retrasos generalizados de 15 a 20~minutos en Metro, no lo era: cada tren se pegaba a la
salida siguiente, la aplicación mostraba «+2~min» o «4~min adelantado» a trenes que iban 18~minutos tarde y dejaba
salidas sin tren como «solo horario». Ese emparejamiento queda solo de reserva, para cuando ningún viaje del feed esté
en el horario (por ejemplo, un horario recién cambiado).

\paragraph{Comprobaciones del Directo} Antes de mostrar nada se comprueba, con los datos de cada consulta:
\begin{itemize}
  \item \emph{Un feed que copia el horario.} Si todas las horas del feed son las del horario al segundo, no es
        directo y sale el horario. Es el caso de los trenes de Euskotren (2\,504 horas, ninguna distinta); los tranvías,
        en el mismo feed, sí traen horas distintas.
  \item \emph{Unas horas que contradicen al GPS.} Si hay al menos 3~vehículos con posición y, en la mayoría de ellos,
        la hora del feed en la parada por la que van se aparta más de 20~minutos de la hora actual
        (\texttt{VEHICLE\_DISAGREEMENT\_SECONDS}), las horas del feed no valen. Es el caso de Álavabus, que las da con una
        hora de más.
  \item \emph{Un viaje que aún no ha salido.} Si no ha salido de su origen y el operador solo repite su hora
        publicada (Renfe, hasta una hora antes, con retraso cero), sale con el horario.
\end{itemize}

\paragraph{Vehículos con Posición y sin Predicción} Su retraso es el que llevan en la parada por la que van (nunca un
adelanto) y se aplica a las paradas que les quedan. No cuentan los que están en la última parada de su viaje ni los que
darían más de 45~minutos (\texttt{MAX\_GPS\_DELAY\_SECONDS}), que suelen ser vehículos parados o que aún hacen el viaje
anterior.

\paragraph{Predicción Vencida} Si un vehículo aún no ha llegado a la parada a la que va y la hora que le daba el
operador ya pasó hace más de 60~s, va al menos ese tiempo más tarde: se suma a esa parada y a las siguientes, y crece
mientras no avance. Con más de una hora de desfase no se extrapola (\texttt{MAX\_STALE\_PREDICTION\_SECONDS}).

\paragraph{Salidas sin Datos y Salidas Pasadas} En los trenes, una salida que ya debería estar en la estación y no está en el
feed va al final como «sin datos del tren». Se miran también las salidas de la última hora, para que un tren muy
retrasado siga en la lista, y se retiran las que ya salieron: las que están en directo, 45~s después de su llegada, y las
programadas, 60~s después.

\paragraph{Vigilancia} El servidor publica para cada red un diagnóstico (viajes del feed, cuántos encuentra en el
horario, cuántos deberían ir en marcha según el horario, si copia el horario o contradice al GPS y la mediana y el
percentil 95 de los retrasos que muestra), y un proceso lo revisa cada hora: si un feed deja de responder, si una red
deja de encontrar sus viajes o si los retrasos no son creíbles (mediana de más de 30~minutos o percentil 95 de más de
60), avisa. Un feed vacío, o uno cuyos viajes no están en el horario, solo cuenta como fallo si según el horario van en
marcha al menos tres viajes de la red; el vacío, además, si la red no es de las que lo publican vacío a ratos aunque
haya servicio (Bilbobus, Lurraldebus y Álavabus). Un feed con vehículos pero ninguno con hora de hoy, como el de Bizkaibus
de madrugada, queda como aviso. Así se detectaron las horas desfasadas de Álavabus, vehículos ya llegados que seguían
contando y trenes de madrugada con $\pm1\,440$~min de retraso.

\subsection{Pasos Para Otra Ciudad}
\label{sec:pasos}

\begin{enumerate}
  \item Comprobar qué publica el operador. Con Trip Updates, basta el emparejamiento de la sección anterior y sus
        comprobaciones. Con SIRI-VM
        (o GTFS-Realtime Vehicle Positions), hace falta el motor de este trabajo.
  \item Construir la base del horario con la fusión de variantes y los alias (sección «\nameref{sec:alias}») y los trazados
        con distancias.
  \item Adaptar la lectura del feed y el emparejamiento, y comprobar con capturas reales que todos los vehículos del
        feed encuentran su viaje.
  \item Poner en marcha la captura periódica y la recalibración diaria. Sin calibrar, el modelo tiene poco con qué
        corregir el horario.
  \item Medir con el protocolo del capítulo «\nameref{cap:evaluacion}» sobre capturas de la propia red antes de mostrar
        nada como «en directo», y volver a medir las constantes de la sección «\nameref{sec:adaptar}».
  \item Comprobar que una implementación propia reproduce los vectores del Apéndice~\ref{ap:vectores}.
\end{enumerate}
